\chapter{Distrito Capital}\label{dept:Distrito Capital}
\mapaparaguai{0}{Distrito Capital}
\vfill\clearpage
\section{Pacote 02: Distrito Capital /\allowbreak 
Asuncion}

\subsection{Leitura executiva}

Distrito Capital funciona como o nó administrativo, portuário e
institucional do país. Para um leitor que busca autossuficiência ou
preparação, o interesse do distrito não está em terra rural, mas em
acesso a serviços, governança, logística, saúde de alta complexidade e
rede de conexões nacionais.

O pacote 02 consolidou uma leitura em seis eixos:

\begin{itemize}
\tightlist
\item
  estrutura urbana e portuária
\item
  energia e tarifas
\item
  saúde, ensino superior e segurança
\item
  mobilidade e corredores logísticos
\item
  vigilância e capacidade institucional
\item
  lacunas operacionais que exigem cautela analítica
\end{itemize}

\subsection{O que o pacote confirma}

\begin{itemize}
\tightlist
\item
  Asunción reúne o principal porto nacional e a administração portuária
  da ANNP, reforçando o papel do Distrito Capital como nó logístico e
  institucional.
\item
  O recorte demográfico e territorial do distrito está disponível no INE
  via Censo 2022, mas as planilhas detalhadas de população ainda não
  foram extraídas neste ciclo.
\item
  A oferta energética é nacional e operada pela ANDE\footnote{Administración Nacional de Electricidade (Administração Nacional de Eletricidade), companhia estatal de energia.}; a referência
  tarifária consolidada no pacote é o Pliego de Tarifas Nro 21, versão
  atualizada de 27-11-2024, já lida para BT/\allowbreak MT no material de apoio.
\item
  A área metropolitana concentra hospitais especializados de alta
  complexidade, incluindo IMT, INCAN, INE\footnote{Instituto Nacional de Estadística (Instituto Nacional de Estatística), órgão oficial de dados demográficos e censitários.}RAM, Hospital General Barrio
  Obrero, Hospital del Trauma, Hospital Materno Infantil San Pablo,
  Hospital Materno Infantil Santísima Trinidad, Hospital Materno
  Infantil Loma Pyta, Hospital Psiquiátrico e Hospital del Indígena.
\item
  O sistema universitário habilitado inclui instituições como
  Universidad Americana, Universidad Autónoma de Asunción e Universidad
  Adventista del Paraguay, entre outras listadas pelo CONES.
\item
  A capital concentra estruturas penitenciárias de referência, com
  destaque para Tacumbú e a transição institucional do Buen Pastor.
\item
  A malha de segurança passa pelo sistema 911, com monitoramento e
  integração de câmeras urbanas em Asunción.
\item
  A conectividade logística da área metropolitana inclui o Aeroporto
  Internacional Silvio Pettirossi, em Luque, como principal terminal
  aéreo do entorno da capital.
\item
  O serviço postal nacional opera com cobertura em Asunción por meio da
  Dinacopa /\allowbreak  Correos del Paraguay.
\item
  O MIC registrou Asunción como principal polo de novos prestadores
  formais do REPSE no 1º quadrimestre de 2025, reforçando o caráter de
  centro comercial e de serviços do distrito.
\item
  O INE\footnote{Instituto Nacional de Estadística (Instituto Nacional de Estatística), órgão oficial de dados demográficos e censitários.} fechou a projeção populacional de 2025 com 464.185 habitantes,
  mediana de idade de 33 anos, fecundidade de 1,68 filho por mulher e
  anos de estudo médios acima de 11 anos.
\item
  A leitura demográfica e educacional saiu da zona de lacuna aberta e
  entrou na camada de dados consolidados.
\end{itemize}

\subsection{Ficha climática de Assunção (NASA POWER,
2001--2020)}

\textbf{Irradiância solar global --- ALLSKY\_SFC\_SW\_DWN (kWh/\allowbreak m²/\allowbreak dia):}
Jan 6.76 \textbar{} Fev 6.20 \textbar{} Mar 5.42 \textbar{} Abr 4.44
\textbar{} Mai 3.36 \textbar{} Jun 2.83 \textbar{} Jul 3.21 \textbar{}
Ago 3.93 \textbar{} Set 4.61 \textbar{} Out 5.49 \textbar{} Nov 6.42
\textbar{} Dez 6.78 Média anual: 4.95 kWh/\allowbreak m²/\allowbreak dia. Potencial solar
superior à média europeia em todos os meses. Mínimo de inverno (2.83 em
junho) ainda viável para sistemas off-grid.

\textbf{Precipitação --- PRECTOTCORR (mm/\allowbreak dia):} Jan 4.06 \textbar{} Fev
4.92 \textbar{} Mar 4.37 \textbar{} Abr 4.97 \textbar{} Mai 4.62
\textbar{} Jun 2.47 \textbar{} Jul 2.03 \textbar{} Ago 1.23 \textbar{}
Set 2.44 \textbar{} Out 5.27 \textbar{} Nov 6.35 \textbar{} Dez 5.78
Média anual: \textasciitilde1.471 mm/\allowbreak ano. Estação chuvosa Out--Mai;
estação seca Jun--Set. Mínimo em agosto (\textasciitilde38 mm/\allowbreak mês) ---
janela estratégica para obras e instalações externas.

\textbf{Ângulo de inclinação para painéis solares:} 25° ao Norte
(anual); 35° no inverno (Jun--Ago); 15° no verão (Nov--Jan). Fonte: NASA
POWER Climatology API (acesso em 2026-03-20).

\subsection{O que ainda falta
consolidar}

\begin{itemize}
\tightlist
\item
  Custo médio formal de moradia e terra por bairro/\allowbreak distrito.
\item
  Base oficial única para postos de combustível, carregamento elétrico,
  gás encanado e inventário comercial.
\item
  IDH e índices sociais padronizados por distrito.
\item
  Índice de violência por 100 mil habitantes em fonte oficial aberta.
\item
  Poluição luminosa: requer proxy satelital (NOAA/\allowbreak VIIRS).
\end{itemize}

\subsection{Síntese editorial}

Asunción não aparece aqui como ``destino de fuga'', mas como
infraestrutura de apoio. Ela importa porque concentra o que um projeto
no interior precisa quando houver dependência de governo, saúde,
educação, carga regulatória ou tráfego internacional.
\clearpage
\newpage
\secaoDiagnostico{Asuncion}{\mapaConteudo{-25.2666}{-57.6333}}{Asunción é uma capital segura para os padrões globais, mas sofre dos
riscos inerentes a qualquer centro de comando nacional: centralização de
alvos e dependência de redes. Ao mesmo tempo, concentra o porto
nacional, a administração pública, hospitais de alta complexidade,
universidades, o sistema 911, o Aeroporto Silvio Pettirossi e a maior
densidade de serviços do país, o que a torna um polo de apoio e não um
destino autônomo.

\subsubsection{Por que sim}

\begin{itemize}
\tightlist
\item
  Melhor infraestrutura de saúde e serviços do país.
\item
  Estabilidade geológica e segurança institucional.
\item
  Rede logística, energética e regulatória mais completa do Paraguai.
\end{itemize}

\subsubsection{Por que não}

\begin{itemize}
\tightlist
\item
  Concentração de alvos militares e governamentais.
\item
  Densidade populacional elevada.
\item
  Alta dependência de infraestrutura centralizada (energia, água,
  logística).
\end{itemize}

\subsubsection{Recomendação
Técnica}

Ideal como base logística e administrativa, desde que combinada com um
refúgio de nível superior em departamentos mais remotos (como Cordillera
ou Itapúa). Referencia atual de score: GSS 6.5 em 2026-03-05. A leitura
do pacote confirma uma capital funcional para comando, saúde e
conectividade, mas não uma base de autossuficiência plena.

\subsubsection{}

\begin{itemize}
\tightlist
\item
  \begin{enumerate}
  \def\labelenumi{\arabic{enumi})}
  \tightlist
  \item
    Dados censitários validados (INE\footnote{Instituto Nacional de Estadística (Instituto Nacional de Estatística), órgão oficial de dados demográficos e censitários.} 2022).
  \end{enumerate}
\item
  \begin{enumerate}
  \def\labelenumi{\arabic{enumi})}
  \setcounter{enumi}{1}
  \tightlist
  \item
    Indicadores de segurança (Ministerio del Interior 2025).
  \end{enumerate}
\item
  \begin{enumerate}
  \def\labelenumi{\arabic{enumi})}
  \setcounter{enumi}{2}
  \tightlist
  \item
    Mapas de inundação e clima (DMH/\allowbreak SEN\footnote{Secretaría de Emergencia Nacional (Secretaria de Emergência Nacional), órgão governamental de resposta a desastres.} 2026).
  \end{enumerate}
\item
  \begin{enumerate}
  \def\labelenumi{\arabic{enumi})}
  \setcounter{enumi}{3}
  \tightlist
  \item
    Tarifário e infraestrutura energética (ANDE\footnote{Administración Nacional de Electricidade (Administração Nacional de Eletricidade), companhia estatal de energia.} 2026).
  \end{enumerate}
\end{itemize}}
\begin{table}[ht!]
\small \begin{tabular}{p{3.0cm}p{0.8cm}p{6.0cm}} \toprule
\textbf{Critério} & \textbf{Nota} & \textbf{Justificativa} \\ \midrule
A - Ameaças Estratégicas & 5.0 & Capital nacional; alvos militares e governamentais presentes \\
B - Risco Social & 5.0 & Alta densidade metropolitana; risco de convergência de refugiados \\
C - Riscos Naturais & 7.5 & Geologia estável; vulnerabilidade hídrica em áreas específicas \\
D - Autossuficiência & 7.0 & Abundância nacional de energia e água; solo fértil periférico \\
E - Institucional & 8.5 & Forte segurança institucional; baixa taxa de criminalidade violenta \\
\midrule \textbf{GSS FINAL} & \textbf{6.5} & \textbf{Moderadamente Seguro (Recomendado para centros de operação com apoio rural).} \\ \bottomrule \end{tabular}
\end{table}
\subsubsection{Vulnerabilidades e
Mitigação}\label{vuln-Asuncion-vulnerabilidades-e-mitigauxe7uxe3o}

\begin{itemize}
\tightlist
\item
  \textbf{Centralização Política:} Alvo óbvio em crises institucionais
  ou conflitos (Mitigação: manter residência secundária em zona
  rural/\allowbreak periférica).
\item
  \textbf{Dependência de Serviços:} Colapso de rede elétrica ou água
  atinge o núcleo urbano rapidamente (Mitigação: geradores solares e
  cisternas).
\item
  \textbf{Risco de Inundação:} Evitar áreas de baixa cota próximas ao
  rio.
\item
  \textbf{Agitação Social:} Centro de manifestações nacionais
  (Mitigação: localização fora do eixo central/\allowbreak administrativo).
\end{itemize}
\subsubsection{Dossiê de Campo}\label{dossie-Asuncion-dossiuxea-de-campo}
\begin{description}[style=nextline,leftmargin=0.5cm,font=\bfseries]
\item[Coordenadas]  25°16'S 57°38'W.
\item[Topografia]  Colinas suaves (geologicamente estável, risco sísmico desprezível). Margeada pelo Rio Paraguai.
\item[Alvos Estratégicos]  Palácio de los López (Governo), Porto Sajonia (Marinha), Base Ñu Guasu (Aeronáutica/\allowbreak Exército), Aeroporto Silvio Pettirossi (Luque), Hub elétrico de transmissão de Itaipu.
\item[Fallout]  Ventos predominantes: Nordeste (verão - Sirocco, quente) e Sul (inverno - Pampero, frio). Baixo risco de fallout global, mas vulnerável se houver ataques a centros de comando nacionais.
\item[Mapa e media local]  conjunto visual de apoio disponível em `MEDIA.md`, com midia institucional local para uso offline e verificação territorial.
\item[População]  464.185 em 2025 (INE\footnote{Instituto Nacional de Estadística (Instituto Nacional de Estatística), órgão oficial de dados demográficos e censitários.}). Censo 2022 Final: 462.241. Área metropolitana (Gran Asunción) com \textasciitilde 2.200.000.
\item[Densidade]  \textasciitilde 3.950 hab/\allowbreak km² (Núcleo urbano denso).
\item[Vias de Acesso]  Hub central do país. Rutas PY01, PY02, PY03. Gargalos em pontes sobre o Rio Paraguai (Remanso, Héroes del Chaco).
\item[Serviços]  Melhor infraestrutura hospitalar e educacional do país. Investimentos maciços em novos centros de saúde.
\item[Custo de Vida]  US\$ 600-1.550/\allowbreak mês (dependendo do padrão). Terrenos urbanos: US\$ 400-650/\allowbreak m².
\item[Saúde de alta complexidade]  IMT, INCAN, INE\footnote{Instituto Nacional de Estadística (Instituto Nacional de Estatística), órgão oficial de dados demográficos e censitários.}RAM, Hospital General Barrio Obrero, Hospital del Trauma, Hospital Materno Infantil San Pablo, Hospital Materno Infantil Santísima Trinidad, Hospital Materno Infantil Loma Pyta, Hospital Psiquiátrico e Hospital del Indígena concentram oferta especializada na capital.
\item[Universidades habilitadas]  Asunción concentra oferta do CONES, com instituições como Universidad Americana, Universidad Autónoma de Asunción e Universidad Adventista del Paraguay.
\item[Conexões logísticas]  Porto de Asunción, sistema 911 e Aeroporto Internacional Silvio Pettirossi reforçam a função de nó administrativo, portuário e metropolitano.
\item[Idade mediana]  33 anos em 2025, segundo projeções do INE\footnote{Instituto Nacional de Estadística (Instituto Nacional de Estatística), órgão oficial de dados demográficos e censitários.}.
\item[Fecundidade]  1,64 filho por mulher em 2022 e 1,68 em 2025 nas projeções do INE\footnote{Instituto Nacional de Estadística (Instituto Nacional de Estatística), órgão oficial de dados demográficos e censitários.}.
\item[Escolaridade]  11,7 anos de estudo em homens e 11,4 em mulheres na faixa de 15+ anos; a cidade concentra o maior capital educacional do país.
\item[Assistência escolar]  os indicadores nacionais da EPHC 2025 mostram 39,4\% de frequência geral e 97,7\% de assistência entre 5 e 9 anos.
\item[Hidrologia]  Risco de Inundação Recorrente. Áreas ribeirinhas (Bañados) sofrem cheias periódicas ligadas ao El Niño.
\item[Clima]  Tropical/\allowbreak Subtropical. Verões com temperaturas >40°C. Tempestades severas com rajadas de vento de até 160 km/\allowbreak h.
\item[Sismicidade]  Muito baixa. Região tectonicamente estável.
\item[Irradiância solar global — ALLSKY\_SFC\_SW\_DWN (kWh/\allowbreak m²/\allowbreak dia)]  Jan 6.76 | Fev 6.20 | Mar 5.42 | Abr 4.44 | Mai 3.36 | Jun 2.83 | Jul 3.21 | Ago 3.93 | Set 4.61 | Out 5.49 | Nov 6.42 | Dez 6.78. Média anual: 4.95 kWh/\allowbreak m²/\allowbreak dia. Mínimo: Jun (2.83) | Máximo: Dez (6.78). Fonte: NASA POWER Climatology API, período 2001-2020 (acesso em 2026-03-20).
\item[Precipitação — PRECTOTCORR (mm/\allowbreak mês)]  Jan 4.06 | Fev 4.92 | Mar 4.37 | Abr 4.97 | Mai 4.62 | Jun 2.47 | Jul 2.03 | Ago 1.23 | Set 2.44 | Out 5.27 | Nov 6.35 | Dez 5.78. Média anual: \textasciitilde 1.471 mm/\allowbreak ano. Estação chuvosa Out-Mai; seca Jun-Set (mínimo em Ago). Fonte: NASA POWER (acesso em 2026-03-20).
\item[Inclinação recomendada para placas solares]  25° voltado para o Norte (anual); 35° no inverno (Jun-Ago); 15° no verão (Nov-Jan). Base: latitude local -25.28°S.
\item[Energia]  Dependência total da rede nacional (Itaipu), mas com a energia mais barata da região (\textasciitilde US\$ 0,05/\allowbreak kWh).
\item[Água]  Abundante (Rio Paraguai), tratamento centralizado pela ESSAP\footnote{Empresa de Servicios Sanitarios del Paraguay (Empresa de Serviços Sanitários do Paraguai), responsável pelo saneamento e água urbana.}.
\item[Solo]  Solo fértil na Região Oriental, mas área urbana saturada. Potencial agrícola nos arredores (Central/\allowbreak Chaco'i).
\item[Autossuficiência]  Baixa no núcleo urbano; alta dependência de suprimentos logísticos nacionais e internacionais.
\item[Internet de alta velocidade]  110 pontos de WiFi libre foram habilitados em Asunción pelo MITIC; 320 novos pontos chegaram ao país em 2023, com cobertura nacional ampliada.
\item[Rede celular]  a tecnologia 5G foi lançada oficialmente no país em dezembro de 2025 e teve ato em Asunción, no Paseo La Galería.
\item[Poços e águas subterrâneas]  não há profundidade média oficial consolidada para poços artesianos em Asunción; a regulação permanece com MADES e Lei 3239/\allowbreak 07.
\item[Combustível e gás]  não foi identificado inventário oficial único por distrito para postos, eletropostos e distribuição de gás encanado.
\item[Serviço público digital]  a municipalidade de Asunción iniciou processo de transformação digital com apoio do MITIC em 2025.
\item[Segurança]  Taxa de homicídios baixa para padrões latinos (7.1/\allowbreak 100k). Criminalidade focada em delitos contra o patrimônio. Áreas nobres e corporativas são consideradas seguras.
\item[Leis Local]  Sede do governo central. Direitos de propriedade bem definidos. Incentivos fiscais (Lei 60/\allowbreak 90).
\item[Delegacias]  Dirección de Policía de Asunción com 23 comisarías jurisdicionais listadas pela Policía Nacional.
\item[Presídios]  Tacumbú concentra 2.151 PPL em censo recente do Ministério de Justicia; o conjunto penitenciário de Asunción também inclui UPIE e Granja Ko'e Pyahu, com mais de 2.200 PPL em visita institucional de 2025.
\item[Correios e entregas]  Correos del Paraguay opera cobertura nacional e atende a capital; não há tabela oficial única de desempenho local por bairro.
\item[Comércio urbano]  o MIC publicou o REPSE do 1º quadrimestre de 2025 com Asunción concentrando 32,2\% dos novos prestadores de serviços do país, o que funciona como proxy oficial de dinamismo comercial formal; o inventário centralizado de cinemas, shoppings e atacados segue sem tabela única.
\item[Escolaridade e ensino superior]  INE\footnote{Instituto Nacional de Estadística (Instituto Nacional de Estatística), órgão oficial de dados demográficos e censitários.} e CONES são as bases centrais para alfabetização, escolaridade e oferta universitária; a extração fina por recorte territorial permanece parcial.
\item[Demografia social]  idade mediana de 33 anos e fecundidade de 1,68 filho por mulher em 2025; o item já saiu da zona de lacuna aberta.
\item[Poluição luminosa]  Bortle 8 (céu urbano), radiância artificial estimada em 100 nW/\allowbreak cm²/\allowbreak sr. Fonte: World Atlas of Artificial Night Sky Brightness.
\end{description}
\paragraph{Solo (SoilGrids 2.0, média ponderada 0--30
cm)}

\begin{longtable}[]{@{}ll@{}}
\toprule\noalign{}
Parâmetro & Valor \\
\midrule\noalign{}
\endhead
\bottomrule\noalign{}
\endlastfoot
pH (H$_2$O) & 5.72 \\
Carbono orgânico (SOC) & 24.33 g/\allowbreak kg \\
Argila & 19.82 \% \\
Areia & 56.72 \% \\
Silte (calc.) & 23.5 \% \\
Densidade aparente & 1.21 g/\allowbreak cm³ \\
\textbf{Aptidão agrícola} & \textbf{Alta} \\
\end{longtable}

Fonte: ISRIC SoilGrids 2.0 via WCS (acesso em 2026-03-21). Coords:
-25.2667°, -57.6333°. Média ponderada camadas 0-5, 5-15, 15-30 cm.

\begin{itemize}
\tightlist
\item
  \textbf{Autossuficiência:} Baixa no núcleo urbano; alta dependência de
  suprimentos logísticos nacionais e internacionais.
\item
  \textbf{Internet de alta velocidade:} 110 pontos de WiFi libre foram
  habilitados em Asunción pelo MITIC; 320 novos pontos chegaram ao país
  em 2023, com cobertura nacional ampliada.
\item
  \textbf{Rede celular:} a tecnologia 5G foi lançada oficialmente no
  país em dezembro de 2025 e teve ato em Asunción, no Paseo La Galería.
\item
  \textbf{Poços e águas subterrâneas:} não há profundidade média oficial
  consolidada para poços artesianos em Asunción; a regulação permanece
  com MADES e Lei 3239/\allowbreak 07.
\item
  \textbf{Combustível e gás:} não foi identificado inventário oficial
  único por distrito para postos, eletropostos e distribuição de gás
  encanado.
\item
  \textbf{Serviço público digital:} a municipalidade de Asunción iniciou
  processo de transformação digital com apoio do MITIC em 2025.
\end{itemize}
